% results.tex (v15, 2026-10-21): five subsections that follow the questions of the Introduction; qualitative reading added after the numbers; external CALM regions removed. Delta = RMSE(first) - RMSE(second) on the same site-years; 95% paired block-bootstrap interval in brackets.
\section{Results}\label{sec:results}

\subsection{Climate input, validation design and the in-region versus out-of-region contrast}
We first ask whether the climate of the observation year matters, how accurate each method is at sites it has not seen, and how much this answer depends on the way the data are split. Replacing the 2011--2024 mean climate input by the input of the observation year reduced the spatial-block RMSE of the recalibrated Stefan model in the Lena Delta ($\Delta$RMSE \ndforcLePo{} cm) and that of the physics-residual model in the Lena Delta (\ndforcLeRo{} cm) and Canada (\ndforcCaRo{} cm), whereas in Alaska, where half of the labels come from one year, the change was not resolved (Fig.~\ref{fig:accuracy}a); beyond accuracy, the year-specific input is what ties each prediction to a stated year. Under spatial-block cross-validation with this input, the physics-residual model reduced the RMSE of the recalibrated Stefan model from \nrmseLeblkPo{} to \nrmseLeblkRo{} cm in the Lena Delta ($\Delta$RMSE \ndrefLeblkRo{} cm) and from \nrmseCablkPo{} to \nrmseCablkRo{} cm in Canada (\ndrefCablkRo{} cm), whereas in Alaska the two models reached \nrmseAkblkPo{} and \nrmseAkblkRo{} cm (\ndrefAkblkRo{} cm; Fig.~\ref{fig:accuracy}b). The scatter of observed against predicted values shows what the learner adds: the recalibrated Stefan model predicts a narrow range of values in every region, because one coefficient applied to the climate of each year cannot separate the sites, whereas the physics-residual model spreads the predictions along the 1:1 line in the Lena Delta and Canada (Supplementary Fig.~S11). Direct ML, which learns the whole response from the labels, reached lower point errors than the physics-residual model when all labels of a region were used (Alaska \nrmseAkblkDz{} versus \nrmseAkblkRo, Lena Delta \nrmseLeblkDz{} versus \nrmseLeblkRo{} and Canada \nrmseCablkDz{} versus \nrmseCablkRo{} cm), but its difference from the recalibrated Stefan model was resolved only in Canada (\ndrefCablkDz{} cm). When the same predictions are scored on 1 km cell-year means instead of site-years, the advantage of direct ML disappears in Alaska (\ncyAkblkDz{} versus \ncyAkblkRo{} cm) and narrows elsewhere (Supplementary Table~S6), which indicates that part of it comes from fitting differences between probe points within cells, which a 1 km map cannot use.

\paperfig{Fig3_accuracy}{1.0}{\textbf{Climate input, validation design and the in-region versus out-of-region contrast.} \textbf{a}, Change in spatial-block RMSE when the 2011--2024 mean climate input is replaced by the input of the observation year (95\% paired block-bootstrap intervals; negative values favour the year-specific input). \textbf{b}, RMSE of three methods under five validation designs, ordered from the least to the most demanding; the colour of a cell is its difference from the recalibrated Stefan model of the same design (blue, lower; red, higher), and in the region hold-out the Stefan model uses the coefficient of the other two regions. \textbf{c}, RMSE of the same methods when a region's own labels are used (spatial-block cross-validation) and when the region is withheld and the models are trained on the other two regions only.}{fig:accuracy}

The ranking of the methods depended on the validation design (Fig.~\ref{fig:accuracy}b, Supplementary Table~S4). Random cross-validation, which places test sites next to training sites, gave the lowest errors to direct ML (\nrmseAkrandDz{} cm in Alaska); spatial-block cross-validation raised the direct-ML error by \noptimAkDz, \noptimLeDz{} and \noptimCaDz{} cm in Alaska, the Lena Delta and Canada, and the year and block-by-year designs raised it further. When an entire region was withheld, the direct-ML error in the Lena Delta rose to \nrmseLeregDz{} cm, against \nrmseLeregPz{} cm for the Stefan model with the coefficient of the other regions ($\Delta$RMSE \ndrefLeregDz{} cm), whereas the physics-residual model stayed within \ndrefLeregRoabs{} cm of the Stefan model there and improved on it in Canada (\ndrefCaregRo{} cm; Fig.~\ref{fig:accuracy}c). The three methods thus behave differently as test sites move away from training sites: the Stefan model is nearly insensitive to the design because one coefficient cannot overfit, direct ML is the most accurate where test and training sites interleave and the least accurate where they do not, and the physics-residual model keeps most of the in-region gain while inheriting the stability of its anchor. Because maps are needed at cells and in years without observations, spatial-block cross-validation is the main design and the region hold-out the test of transferability.

\subsection{Conditions of the physics-residual gain}
We next ask what information the learner needs and where the correction acts. The gain depended on the covariates available to the residual learner (Fig.~\ref{fig:conditions}a): with climate covariates only, the residual model did not improve on the Stefan model in Alaska (\ndfeatAkFclim{} cm) or Canada (\ndfeatCaFclim{} cm), and adding terrain and soil covariates produced the gains reported above; MODIS vegetation and snow timing or high-resolution soil and land-cover classes changed the result by less than 0.5 cm in the Lena Delta and Canada, and the latter added a small gain in Alaska (\ndfeatAkFhires{} cm). Across regional terciles of eight covariates, the gain occurred in most conditions of the Lena Delta (\ngcLesig{} of \ngcLecells{} terciles with intervals excluding zero) and Canada (\ngcCasig{} of \ngcCacells) and increased with the error of the Stefan model (correlation across terciles \ngcLecorr{} in the Lena Delta and \ngcCacorr{} in Canada; Fig.~\ref{fig:conditions}b), reaching \ndtddLewarm{} cm in the warmest TDD tercile of the Lena Delta. In Alaska no tercile changed by more than \ngcAkmaxabs{} cm; the absence of a gain is uniform and consistent with two features of the Alaska data: the regional coefficient already represents its sites well (Stefan RMSE \nrmseAkblkPo{} cm, the lowest of the three regions relative to the range of the observations), and the remaining error is dominated by the spread between probe points within 1 km cells (pooled within-cell standard deviation \ncyWithinSdAk{} cm), which no 1 km covariate can resolve. Terrain variables accounted for \nshapLeResTerrain\% of the mean absolute SHAP value of the residual model in the Lena Delta and soil variables for \nshapCaResSoil\% in Canada (Supplementary Fig.~S3). Placing the physics elsewhere, as pseudo-labels (which reproduced the Stefan model, Lena Delta \ndrefLeblkDo{} cm) or as an input feature (which behaved like direct ML), and replacing CatBoost by four other learners did not change the sign of the gain in the Lena Delta and Canada (Supplementary Fig.~S2).

\paperfig{Fig4_conditions}{1.0}{\textbf{Conditions of the physics-residual gain and transfer to a region with few labels.} \textbf{a}, Difference in spatial-block RMSE between the physics-residual model and the recalibrated Stefan model for six nested covariate sets. \textbf{b}, The same difference within regional terciles of eight covariates; outlined cells have intervals excluding zero, hatched cells fewer than 30 site-years or five blocks. \textbf{c}, RMSE in each region for models trained on the other two regions plus $n$ labels of the target region; a value beyond the axis is printed. Intervals in a and b are 95\% paired block-bootstrap intervals; negative differences favour the physics-residual model.}{fig:conditions}

\subsection{Number of labels within a region and transfer to a new region}
The third question is how many local observations the learner needs. With 25 training labels drawn from the site-years of a region, the physics-residual model was at least as accurate as a Stefan model recalibrated on the same labels (Alaska \nllAkRotwfive, Lena Delta \nllLeRotwfive{} and Canada \nllCaRotwfive{} cm), and with 100 labels its gain was resolved in all three regions (\nllAkRohundred, \nllLeRohundred{} and \nllCaRohundred{} cm; Supplementary Fig.~S4). Direct ML trained on 25 labels had a higher point error than the Stefan model in Alaska (\nllAkDztwfive{} cm) and Canada (\nllCaDztwfive{} cm) and overtook the physics-residual model only with several hundred labels in Alaska and the Lena Delta. The two curves do not contradict each other: with thousands of labels direct ML fits them more closely, but the physics-residual model reaches almost the same accuracy with a fraction of the labels and never falls below its anchor. Other model structures: Supplementary Fig.~S4.

In the transfer experiment, each region was predicted by models trained on the other two regions and $n$ labels of the target region (Fig.~\ref{fig:conditions}c; $n \leq 500$, drawn from the training folds of the target region). Without target labels, direct ML reached \ntlrmseLeDszero{} cm in the Lena Delta, against \ntlrmseLePzzero{} cm for the source-coefficient Stefan model and \ntlrmseLeRszero{} cm for the physics-residual model; 25 target labels reduced the physics-residual error to \ntlrmseLeRstwfive{} cm and the direct-ML error to \ntlrmseLeDstwfive{} cm. In Alaska the physics-residual model gave the lowest error at every $n$ (\ntlrmseAkRszero{} cm without target labels). In Canada, direct ML trained on the source regions was the most accurate (\ntlrmseCaDszero{} cm): recalibrating the Stefan coefficient on 25 Canadian labels increased its error from \ntlrmseCaPzzero{} to \ntlrmseCaPntwfive{} cm, and the physics-residual model, which uses this anchor, followed (\ntlrmseCaRszero{} to \ntlrmseCaRstwfive{} cm). We attribute this to the wide range of thaw depths among the Canadian sites (observed means of \nsrcCaABoVEObs{} cm at the probe transects and \nsrcCaCALMObs{} cm at the CALM grids), which a coefficient fitted to 25 labels represents poorly.

\subsection{Comparison with published products in the same cells and years}
We then ask how the published products perform where our labels are, scored in the same way as our models. On the \nccLeN{} common site-years of the Lena Delta (2010--2020; site-years at which Wei et al., Liu et al. and ESA CCI v4 all provide a value for the observation year), the physics-residual model reached an RMSE of \nccLeRoall{} cm, against \nccLeCcifiveall{} cm for ESA CCI v5, \nccLeAaltoall{} cm for Aalto et al., \nccLeCcifourall{} cm for ESA CCI v4, \nccLeLiuall{} cm for Liu et al. and \nccLeWeiall{} cm for Wei et al. (Fig.~\ref{fig:products}a, Supplementary Table~S5). On the Lena Delta site-years more than 25 km from the training sites of Wei et al. (\nnprodLeweimask{} with a value of Wei et al., \nnprodLeccifivemask{} with a value of ESA CCI v5), the paired differences were \ndprodLeweiRomask{} cm relative to Wei et al. and \ndprodLeccifiveRomask{} cm relative to ESA CCI v5. In Alaska (\nccAkRoall{} versus \nccAkWeiall{} and \nccAkLiuall{} cm for Wei et al. and Liu et al.) and Canada (\nccCaRoall{} versus \nccCaWeiall{} and \nccCaLiuall{} cm) the model and the two ML products were on par.

The scatter of observed against predicted values shows why the regions differ (Fig.~\ref{fig:products}b). In the Lena Delta the two ML products assign most sites values between about 40 and 70 cm regardless of the observed ALT: the correlation between observed and predicted values was \nrLeWei{} for Wei et al. and \nrLeLiu{} for Liu et al., against \nrLeRo{} for the physics-residual model, and their mean bias was \nbiasLeAllWei{} and \nbiasLeAllLiu{} cm, against \nbiasLeAllRo{} cm, with the largest overestimation at the thin active layers of the northern delta (\nbiasLeNorthWei{} and \nbiasLeNorthLiu{} cm north of 72$^\circ$N). The error distributions (Fig.~\ref{fig:products}c) separate this systematic bias from the spread: in the Lena Delta the products are shifted by about 20 cm with a spread similar to that of our model, whereas in Alaska and Canada all series are centred near zero and differ mainly in spread. The regional contrast follows the training data of the products. Alaska label sites lie at a median distance of \ndistweiAkmed{} km from the training sites of Wei et al., and Liu et al. trained their model on the ABoVE active-layer compilation, an earlier release of the Alaska and Canada observations used here, so that their values at these sites are not independent of our labels; the Lena Delta sites lie at a median distance of \ndistweiLemed{} km from the training sites of Wei et al. (Supplementary Fig.~S6). A model trained without any Lena Delta label was already more accurate there than Wei et al. on the same site-years (\ntlvprmseLeRsWeizero{} versus \ntlvprmseLeWeizero{} cm; $\Delta$RMSE \ntlvpLeRsWeizero{} cm), so that the advantage in this region does not depend on local labels alone. Year by year, the physics-residual model had a lower error than Wei et al. in all \nywLen{} Lena Delta years with at least 30 site-years and in \nywAkRoWeiwin{} of \nywAkRoWein{} Alaska years (Supplementary Fig.~S5).

\paperfig{Fig5_products}{1.0}{\textbf{Comparison with published products on the same sites and years.} \textbf{a}, RMSE of the recalibrated Stefan model, the physics-residual model and five products on the common site-years (2010--2020), for all sites (filled) and for sites more than 25 km from the training sites of Wei et al. (open); values above 45 cm are drawn on the compressed segment. \textbf{b}, Observed against predicted ALT on the common site-years for the physics-residual model, Wei et al. and Liu et al., with the RMSE, mean bias and correlation. \textbf{c}, Distribution of predicted minus observed ALT on the same site-years (violins clipped at $\pm$80 cm; dot, median; bar, interquartile range). Our models use spatial-block out-of-fold predictions throughout.}{fig:products}

Maps of the year with the most observed cells in each region (2013 in Alaska and Canada, 2016 in the Lena Delta) show these differences in space (Fig.~\ref{fig:prodmaps}). At the \nfeAkN{} Alaska cells, the RMSE of the three maps was similar (\nfeAkOof{}, \nfeAkWei{} and \nfeAkLiu{} cm). The correlation with the observations was lower for our model (\nfeAkROof) than for the products (\nfeAkRWei{} and \nfeAkRLiu), but \nfeAkNNear{} of the cells lie within 25 km of a training site of Wei et al.; at the \nfeAkNFar{} cells beyond that distance the three maps had similar errors (\nfeAkFarOof{}, \nfeAkFarWei{} and \nfeAkFarLiu{} cm; $r$ \nfeAkFarROof{}, \nfeAkFarRWei{} and \nfeAkFarRLiu). The largest differences between the maps occur above 600 m in mountain terrain; between 68 and 70$^\circ$N the products are deeper than our map there by \nmapAkFootMtnWei{} and \nmapAkFootMtnLiu{} cm. In the Lena Delta (\nfeLeN{} cells observed in 2016) both products were deeper than the observations (mean bias \nfeLeBiasWei{} and \nfeLeBiasLiu{} cm; RMSE \nfeLeWei{} and \nfeLeLiu{} cm against \nfeLeOof{} cm). In Canada (NW), the RMSE at the \nfeCaN{} observed cells was \nfeCaOof{} cm for our model, \nfeCaWei{} cm for Wei et al. (mean bias \nfeCaBiasWei{} cm) and \nfeCaLiu{} cm for Liu et al., whose training data include the earlier release of these observations; at the \nfeCaNFar{} cells beyond 25 km from a training site of Wei et al. the values were \nfeCaFarOof{}, \nfeCaFarWei{} and \nfeCaFarLiu{} cm.

\paperfig{Fig6_product_maps}{1.0}{\textbf{ALT maps from the physics-residual model and two published products for the year with the most observed cells in each region, with the observed values.} \textbf{a--c}, Alaska 2013: physics-residual model, Wei et al. and Liu et al.; the product maps are shown on the cells of our map, and the observed cells are drawn as circles filled with the observed ALT on the same colour scale. \textbf{d}, Observed against predicted ALT at the \nfeAkN{} observed cells with values in all three maps (1:1 line); filled markers are cells within 25 km of a training site of Wei et al. and open markers cells beyond, and the RMSE is given below the panel for all cells and for the open cells; our predictions are spatial-block out-of-fold values. \textbf{e--h}, The same for the Lena Delta 2016 (\nfeLeN{} cells). \textbf{i--l}, The same for Canada (NW) 2013 (\nfeCaN{} cells). Each region has one colour range from the 5th--95th percentiles of the three maps inside the area of applicability; lighter colours mark cells outside it.}{fig:prodmaps}

The residual learner can also be anchored to a product (equation~(\ref{eq:prod})): with the Lena Delta labels the Stefan-anchored model remained more accurate than the product-anchored model (\npadLeRoRp{} cm; Supplementary Fig.~S7), and in Alaska and Canada the two anchors were not distinguishable.

\subsection{Year-labelled 1 km ALT maps, 2011--2024}
Finally, the physics-residual models fitted to all site-years of each region were applied with the ERA5-Land climate of each year from 2011 to 2024 (Fig.~\ref{fig:maps}); the maps cover \nncellmapAk{} cells in Alaska, \nncellmapLe{} in the Lena Delta and \nncellmapCa{} in Canada (NW). Relative to the 2011--2024 mean, the regional mean anomaly ranged from \nanminAk{} cm in \nanminyrAk{} to \nanmaxAk{} cm in \nanmaxyrAk{} in Alaska, from \nanminLe{} cm in \nanminyrLe{} to \nanmaxLe{} cm in \nanmaxyrLe{} in the Lena Delta and from \nanminCa{} cm in \nanminyrCa{} to \nanmaxCa{} cm in \nanmaxyrCa{} in Canada (Fig.~\ref{fig:maps}f). These anomalies are the response of the model to the climate of each year: the regional mean anomaly follows the regional mean TDD to 1 September with a correlation of \ntddanomAk{}, \ntddanomLe{} and \ntddanomCa{} in Alaska, the Lena Delta and Canada (Fig.~\ref{fig:maps}d), so that the deepest regional thaw falls in the warmest thaw seasons (\ntddmaxAk{}, \ntddmaxLe{} and \ntddmaxCa) and the shallowest in the coldest (\ntddminAk, \ntddminLe{} and \ntddminCa); the residual learner modulates this response spatially, so that the anomaly of a warm year is largest in the interior lowlands of Alaska and smallest on the North Slope (Fig.~\ref{fig:maps}f). A comparison with the yearly means of the observed cells is limited by the changing set of sites measured in each campaign, so Supplementary Fig.~S10 compares the mapped and observed series at long-record sites instead. The linear trend of the yearly maps averaged \ntrendAk{}, \ntrendLe{} and \ntrendCa{} cm per decade over the Alaska, Lena Delta and Canada cells, with positive values in \ntrendposAk\%, \ntrendposLe\% and \ntrendposCa\% of the cells (Fig.~\ref{fig:maps}e); it describes the ERA5-Land climate of the period as transmitted by the model and is not an observed trend. The area of applicability covers \naoaAk\% of the Alaska cells, \naoaLe\% of the Lena Delta cells and \naoaCa\% of the Canada cells; cells outside it lie mainly in mountains and at the southern margin of the permafrost zone (Supplementary Fig.~S8). The half-width of the 90\% prediction interval is \nhwAk{}, \nhwLe{} and \nhwCa{} cm. Supplementary Fig.~S9 shows the residual correction.

\paperfig{Fig7_maps}{1.0}{\textbf{Year-labelled 1 km ALT maps, 2011--2024.} \textbf{a--c}, Mean ALT of 2011--2024 from the yearly maps of Alaska, the Lena Delta and Canada (NW), with windows enlarged at the right (1 km cells) and the observed cells filled with their observed mean ALT; lighter colours mark cells outside the area of applicability. \textbf{d}, Regional mean ALT anomaly of the yearly maps against the regional mean TDD anomaly of the same year (one point per year, least-squares line; the warmest and coldest thaw seasons are labelled). \textbf{e}, Linear trend of the yearly maps (least squares per cell), which describes the modelled response to the climate of each year. \textbf{f}, ALT anomaly of each year relative to the 2011--2024 mean in Alaska.}{fig:maps}
